%==============================================================================
%  CHAPITRE 4 — IMPLÉMENTATION DU PROTOTYPE
%  Contenu SEUL (à inclure via \input / \include dans le document principal).
%------------------------------------------------------------------------------
%  Paquets à avoir dans le préambule du document principal :
%    \usepackage{tikz}
%    \usetikzlibrary{arrows.meta,positioning,shapes.geometric,calc,fit,backgrounds}
%    \usepackage{float}
%    \usepackage{booktabs}
%    \usepackage{longtable}   % ou ltablex
%    \usepackage{enumitem}
%    \usepackage[backend=biber,style=numeric-comp,sorting=nyt,maxbibnames=99]{biblatex}
%  Références citées dans ce chapitre (présentes dans references.bib) :
%    hevner2004design, gupta2016daisee, busso2008iemocap, touvron2023llama,
%    lewis2020retrieval, baevski2020wav2vec2, pedregosa2011scikit.
%==============================================================================
\titreduchapitre{Implémentation du prototype}{chap:implementation}

\par Le chapitre précédent a défini \emph{ce que} le système doit faire : percevoir
un état d'engagement à partir du visage et de la voix, décider s'il faut
intervenir, puis générer un contenu pédagogique adapté. Ce chapitre expose
\emph{comment} ces choix de conception ont été réalisés, sous les contraintes
concrètes d'un travail de master. Nous présentons d'abord l'infrastructure
technique et les principes d'architecture qui structurent le code
(section~\ref{sec:infrastructure}), puis le traitement des données
(section~\ref{sec:impl-donnees}), l'implémentation du modèle d'adaptation
(section~\ref{sec:impl-modele}) et celle du module de génération
(section~\ref{sec:impl-generation}). Nous décrivons enfin le prototype de
simulation qui rend le dispositif observable (section~\ref{sec:impl-prototype}) et
la génération du corpus de contenu adaptatif qui en constitue le livrable
(section~\ref{sec:impl-corpus}). Conformément à la démarche de recherche orientée
conception \cite{hevner2004design}, l'artefact logiciel n'est pas ici un simple
support de démonstration : il est l'objet même dont les propriétés seront évaluées
au chapitre~\ref{chap:validation}.

%==============================================================================
\section{Infrastructure technique}
\label{sec:infrastructure}

Quatre contraintes ont orienté les choix techniques exposés dans cette section. Le
système a été développé et exécuté sur une machine de travail ordinaire,
\textbf{sans accélérateur graphique dédié} : toute inférence s'effectue sur
processeur, ce qui exclut l'entraînement de réseaux profonds de bout en bout. Le
prototype devait fonctionner \textbf{sans budget d'infrastructure}, d'où le recours
à des bibliothèques libres exécutées localement et à l'offre gratuite d'un service
de génération distant. Les résultats du chapitre~\ref{chap:validation} devant être
\textbf{reproductibles}, les modèles entraînés, corpus dérivés et métriques sont
des fichiers déposés à des emplacements fixes et tout tirage aléatoire est
initialisé par une graine explicite. Enfin, l'exigence de \textbf{transparence}
posée au chapitre~\ref{chap:conception} impose que chaque étape laisse une trace
lisible plutôt qu'un simple résultat final.

\subsection{Stack technologique}
\label{subsec:stack}

Le prototype n'emploie que deux langages. \textsc{Python}~3.11 couvre
l'intégralité du serveur applicatif \emph{et} des scripts hors ligne (entraînement,
évaluation, construction de corpus) : ce choix garantit que les modules de
perception, de fusion et de décision utilisés par l'interface sont littéralement
les mêmes objets que ceux exécutés par les scripts de validation, et non des
réimplémentations. \textsc{JavaScript} (\textsc{React}) couvre l'interface cliente.
Le tableau~\ref{tab:stack-outils} récapitule la pile retenue, couche par couche.

\begin{table}[H]
\centering
\small
\renewcommand{\arraystretch}{1.15}
\begin{tabular}{|>{\bfseries}l|>{\raggedright\arraybackslash}p{4.3cm}|>{\raggedright\arraybackslash}p{6.1cm}|}
\hline
\rowcolor{gray!20}
Couche & Technologies & Rôle dans le prototype \\
\hline
Exécution & Python 3.11 &
Langage unique du serveur et des scripts hors ligne. \\
\hline
Service web & FastAPI, Uvicorn, Pydantic &
Exposition de l'API REST/JSON et validation des requêtes. \\
\hline
Vision & OpenCV, Transformers, PyTorch &
Décodage des clips, détection du visage, classification des expressions. \\
\hline
Audio & librosa, soundfile, Transformers &
Lecture du signal, extraction du segment annoté, reconnaissance vocale. \\
\hline
Apprentissage & scikit-learn \cite{pedregosa2011scikit}, joblib &
Mapping appris visage $\rightarrow$ engagement et sérialisation des modèles. \\
\hline
Génération & SDK Groq, \textsc{LLaMA}~3.1 \cite{touvron2023llama,dubey2024llama3} &
Production du contenu pédagogique adaptatif. \\
\hline
Récupération & ChromaDB, sentence-transformers (\textsc{MiniLM}) \cite{wang2020minilm} &
Base vectorielle et encodage des documents pour le RAG
\cite{lewis2020retrieval,gao2023retrieval}. \\
\hline
Persistance & SQLite, fichiers JSONL et CSV &
Sessions, journal des décisions, référentiels de scénarios. \\
\hline
Client et tests & React~19, jsPDF, pytest, React Testing Library &
Interface de simulation, export du contenu, tests unitaires et d'API. \\
\hline
\end{tabular}
\caption{Pile technologique du prototype \emolearn{}, par couche fonctionnelle.}
\label{tab:stack-outils}
\end{table}

\paragraph{Réutilisation de modèles pré-entraînés.} La contrainte matérielle a une
conséquence directe sur la perception : aucun réseau profond n'est entraîné depuis
l'initialisation. Des modèles publics sont mobilisés tels quels et seule la
\emph{dernière articulation} est apprise localement — celle qui relie leurs sorties
à la taxonomie d'engagement de \textsc{DAiSEE} \cite{gupta2016daisee}. Le visage
est traité par une cascade de détection OpenCV puis par un classifieur
d'expressions \textsc{Hugging Face} \cite{wolf2020transformers} à sept émotions, tous deux \textbf{gelés}, dont
le vecteur de scores alimente le mapping appris (section~\ref{sec:impl-modele}) ;
la voix repose sur un modèle \emph{wav2vec2} spécialisé en émotions
\cite{baevski2020wav2vec2} dont l'encodeur est gelé et la seule tête de
classification \textbf{ré-entraînée} vers les quatre états \textsc{DAiSEE} à partir
d'\textsc{IEMOCAP} \cite{busso2008iemocap} ; le RAG vectorise ses documents avec un
encodeur multilingue de phrases ; la génération est confiée à un modèle de langue
distant, sans réglage fin, avec une température basse et un budget de sortie borné.

\paragraph{Conséquences de l'exécution sur processeur.} Travailler sans
accélérateur impose des arbitrages explicites. L'analyse faciale d'un clip
n'examine pas toutes ses images mais un \textbf{échantillon de cinq images}
réparties sur sa durée : en deçà, l'estimation devient sensible à une image mal
cadrée ; au-delà, le coût croît linéairement sans gain observé. La génération de
texte est déportée vers un service distant, ce qui évite d'héberger localement un
modèle de plusieurs milliards de paramètres, et les modèles lourds ne sont chargés
qu'à leur première sollicitation effective. Les dépendances sont figées dans un
fichier de spécification unique, l'unique secret — la clé d'accès au service de
génération — est lu depuis un fichier d'environnement non versionné, et les
emplacements sont normalisés : \texttt{data/} (corpus et référentiels),
\texttt{models/} (modèles entraînés), \texttt{results/} (métriques),
\texttt{logs/} (journaux).

\subsection{Choix d'architecture}
\label{subsec:choix-archi}

L'architecture fonctionnelle en trois couches a été posée au
chapitre~\ref{chap:conception}. Il s'agit ici de la traduire en une organisation
logicielle, c'est-à-dire de répondre à des questions d'un autre ordre : quels
processus s'exécutent, comment le code est découpé, où sont stockés les artefacts,
et que se passe-t-il lorsqu'un composant est indisponible
(figure~\ref{fig:archi-technique}).

\begin{figure}[H]
\centering
\begin{tikzpicture}[
  font=\footnotesize,
  box/.style={rectangle, rounded corners, draw=black, thick, align=center,
              text width=3.9cm, minimum height=1.0cm, fill=white},
  ext/.style={rectangle, rounded corners, draw=black, thick, align=center,
              text width=3.0cm, minimum height=1.0cm, fill=black!8},
  store/.style={rectangle, draw=black!60, dashed, align=center,
                text width=10.6cm, minimum height=0.9cm},
  arr/.style={-{Stealth[length=2.6mm]}, thick},
  barr/.style={{Stealth[length=2.6mm]}-{Stealth[length=2.6mm]}, thick}
]
\node[box] (ui) {Navigateur\\interface React};
\node[box, below=1.0cm of ui] (api)
  {Serveur Uvicorn\\API FastAPI (JSON)};
\node[box, below=0.9cm of api] (pipe)
  {Orchestrateur \texttt{pipeline}\\(analyse d'un clip)};
\node[box, below=0.95cm of pipe, xshift=-5.0cm] (perc)
  {Perception\\FER · SER · fusion};
\node[box, right=0.5cm of perc] (gen)
  {Décision\\et génération};
\node[ext, right=1.2cm of gen] (groq) {Service LLM\\distant (Groq)};
\node[store, below=3.0cm of pipe] (art)
  {Artefacts locaux : \texttt{data/} (corpus, scénarios) ·
   \texttt{models/} (FER, SER) · \texttt{data/rag\_db/} (base vectorielle) ·
   \texttt{emolearn.db} (sessions) · \texttt{logs/decisions.jsonl} (journal)};

\draw[barr] (ui) -- node[right, xshift=1mm] {HTTP / JSON} (api);
\draw[arr] (api) -- (pipe);
\draw[arr] (pipe) -- (perc);
\draw[arr] (pipe) -- (gen);
\draw[barr] (gen) -- (groq);
\draw[barr] (perc) -- (art);
\draw[barr] (gen) -- (art);
\end{tikzpicture}
\caption{Vue technique du prototype : un client web, un processus serveur unique,
un service de génération distant et des artefacts persistants locaux (en
pointillés).}
\label{fig:archi-technique}
\end{figure}

\paragraph{Client--serveur et découpage par responsabilité.} Le client et le
serveur sont deux applications distinctes communiquant par des échanges JSON sur
HTTP, et l'API ne conserve \textbf{aucune session en mémoire} : chaque requête
porte son contexte complet (scénario, concept, module, niveau de l'apprenant,
moment du cours) et l'historique nécessaire aux règles temporelles est relu depuis
la couche de persistance. Le serveur peut ainsi être redémarré sans perdre les
parcours simulés, et une analyse peut être rejouée depuis un script sans passer par
l'interface. Le code lui-même n'est pas organisé par couche technique mais par
\emph{responsabilité scientifique}, chaque module étant autonome, testable
isolément et substituable : la perception réunit la résolution et le décodage des
clips (\texttt{daisee\_video}), la chaîne faciale (\texttt{fer\_offline},
\texttt{fer\_classifier}, \texttt{fer\_module}), la chaîne vocale
(\texttt{iemocap\_parser}, \texttt{iemocap\_audio}, \texttt{ser\_module}) et la
combinaison des modalités (\texttt{fusion\_module}) ; la décision repose sur
\texttt{temporal\_rules} et \texttt{decision\_engine} ; la génération sur
\texttt{prompts/templates}, \texttt{rag\_module}, \texttt{llm\_module} et
\texttt{chat\_module} ; l'infrastructure sur \texttt{pipeline} (orchestration et
journalisation), \texttt{database} (sessions) et \texttt{main} (points d'entrée de
l'API).

\paragraph{Un orchestrateur unique, partagé avec le banc d'essai.} La séquence
complète d'analyse d'un clip est écrite \emph{une seule fois}, dans le module
d'orchestration : l'API n'y ajoute que la validation de la requête et la mise en
forme de la réponse. Les scripts hors ligne — évaluation de la perception faciale,
construction du sous-corpus audiovisuel aligné, analyse de sensibilité — importent
les \emph{mêmes} modules de perception et de fusion. Ce point conditionne la
validité du chapitre~\ref{chap:validation} : aucun écart de code ne peut
s'installer entre le système mesuré et le système utilisé.

\paragraph{Commutation des variantes par configuration.} Le
chapitre~\ref{chap:conception} a conservé, pour chaque modalité, une référence
explicable et une variante apprise, ainsi que deux stratégies de résolution de
conflit. Pour que ces alternatives restent comparables, elles sont sélectionnables
par variables d'environnement, sans modification du code :
\texttt{EMOLEARN\_\allowbreak FER\_\allowbreak MAPPING} (\emph{auto} par défaut,
\emph{learned} ou \emph{heuristique}, le mode \emph{auto} ne retenant le modèle
appris que s'il dépasse la référence sur les deux métriques de contrôle),
\texttt{EMOLEARN\_\allowbreak SER\_\allowbreak METHODE} (\emph{learned} ou
\emph{heuristique}) et \texttt{EMOLEARN\_\allowbreak FUSION\_\allowbreak
STRATEGIE} (\emph{priorité pédagogique} ou \emph{confiance}). Les conditions expérimentales du
chapitre~\ref{chap:validation} sont donc des combinaisons de ces variables, et
chaque mesure est rejouable par une commande explicite.

\paragraph{Robustesse d'exécution.} Plusieurs ressources de la chaîne peuvent faire
défaut au moment de l'exécution. Plutôt que d'interrompre le traitement,
l'implémentation
prévoit un repli à chaque point de fragilité : un modèle appris absent ramène au
mapping heuristique de référence, un échec de la récupération documentaire laisse
la génération se poursuivre sans ancrage, une limite de débit du service distant
déclenche des tentatives espacées et l'absence de modalité vocale conduit à une
décision unimodale. Le principe commun est qu'une dégradation
n'est jamais silencieuse : elle est déclarée dans la réponse et dans le journal,
faute de quoi une performance mesurée pourrait être attribuée à une configuration
qui n'était pas celle réellement exécutée.

\paragraph{Traçabilité, apprentissage hors ligne et périmètre assumé.} La
persistance distingue deux besoins : un \textbf{journal de décision}, en ajout
seul, enregistre pour chaque clip l'état fusionné, le type de conflit, le patron
retenu et les indicateurs de génération — c'est la matière première des analyses du
chapitre~\ref{chap:validation} — tandis qu'une \textbf{base de sessions} conserve le
déroulé pédagogique consultable dans l'interface. L'entraînement, lui, reste
\textbf{hors du serveur} : les scripts d'apprentissage produisent des artefacts sur
disque que les modules de perception se contentent de charger, si bien qu'un modèle
ne peut pas changer silencieusement entre deux exécutions. Cette infrastructure
demeure toutefois celle d'un \textbf{prototype de simulation} : base mono-fichier
adaptée à un utilisateur unique, absence d'authentification et de cloisonnement des
accès, aucune conteneurisation. Ces limites correspondent au périmètre de recherche
défini au chapitre~\ref{chap:introduction} et sont discutées au
chapitre~\ref{chap:conclusion} ; ce qui devait être solide ici n'est pas
l'exploitabilité industrielle, mais la reproductibilité et la traçabilité de la
chaîne scientifique.

%==============================================================================
\section{Traitement et préparation des données}
\label{sec:impl-donnees}

Le chapitre~\ref{chap:conception} a justifié le recours à deux corpus
complémentaires : \textsc{DAiSEE} \cite{gupta2016daisee} fournit le signal facial
annoté en états d'engagement, \textsc{IEMOCAP} \cite{busso2008iemocap} fournit le
signal vocal annoté en émotions. Aucun des deux n'est utilisable en l'état par le
système : le premier est distribué comme une arborescence de clips vidéo
accompagnée de fichiers d'annotation, le second comme une collection de dialogues
continus dont les segments pertinents doivent être localisés. Cette section décrit
les traitements qui transforment ces deux ressources brutes en artefacts
exploitables — un référentiel de scénarios, un index d'énoncés vocaux et des jeux de
données d'entraînement — ainsi que le découpage temporel retenu pour la simulation
(figure~\ref{fig:pipeline-donnees}).

\begin{figure}[H]
\centering
\begin{tikzpicture}[
  font=\footnotesize,
  src/.style={rectangle, rounded corners, draw=black, thick, align=center,
              text width=2.9cm, minimum height=0.9cm, fill=black!8},
  op/.style={rectangle, draw=black, thick, align=center,
             text width=3.4cm, minimum height=0.9cm, fill=white},
  res/.style={rectangle, rounded corners, draw=black, thick, align=center,
              text width=3.6cm, minimum height=0.9cm, fill=white},
  arr/.style={-{Stealth[length=2.4mm]}, thick}
]
\node[src] (dlab) {\textsc{DAiSEE}\\annotations};
\node[src, below=0.6cm of dlab] (dvid) {\textsc{DAiSEE}\\clips vidéo};
\node[src, below=0.95cm of dvid] (iemo) {\textsc{IEMOCAP}\\dialogues};

\node[op, right=1.0cm of dlab] (split) {Rééchantillonnage\\70 / 15 / 15};
\node[op, right=1.0cm of dvid] (faces) {Détection du visage\\5 images par clip};
\node[op, right=1.0cm of iemo] (parse) {Analyse des fichiers\\d'annotation};

\node[res, right=1.0cm of split] (scen) {\texttt{scenarios\_50.csv}\\50 séquences};
\node[res, right=1.0cm of faces] (feat) {Cache de traits FER\\(entraînement)};
\node[res, right=1.0cm of parse] (utt) {\texttt{iemocap\_parsed.csv}\\10\,039 énoncés};

\draw[arr] (dlab) -- (split);
\draw[arr] (dvid) -- (faces);
\draw[arr] (iemo) -- (parse);
\draw[arr] (split) -- (scen);
\draw[arr] (faces) -- (feat);
\draw[arr] (parse) -- (utt);
\end{tikzpicture}
\caption{Chaîne de préparation des données : les deux corpus sources (à gauche)
sont transformés hors ligne en trois artefacts versionnés (à droite), qui
alimentent ensuite l'entraînement et la simulation.}
\label{fig:pipeline-donnees}
\end{figure}

\subsection{Extraction et nettoyage du dataset DAiSEE}
\label{subsec:impl-daisee}

\paragraph{Reconstitution d'un référentiel unique.} Les annotations de
\textsc{DAiSEE} sont fournies en trois fichiers correspondant au découpage
d'origine des auteurs (5\,358 clips d'apprentissage, 1\,429 de validation et
1\,784 de test, soit 8\,571 clips annotés). Chaque ligne associe un identifiant de
clip à quatre entiers compris entre 0 et 3 pour l'ennui, l'engagement, la confusion
et la frustration. Les trois fichiers ont été concaténés puis redécoupés en un
partage \textbf{70 / 15 / 15} — 5\,999 clips de conception, 1\,285 de validation et
1\,287 de test — après un mélange initialisé par la graine 42. Ce redécoupage
répond à une exigence méthodologique : le sous-ensemble servant à concevoir les
règles et à choisir les scénarios doit être disjoint de celui qui servira à mesurer
les performances au chapitre~\ref{chap:validation}. Les noms de colonnes sont
normalisés à la lecture, car les fichiers d'origine contiennent des espaces
parasites qui font échouer un accès direct par nom.

\paragraph{Résolution des clips vers les fichiers vidéo.} L'arborescence des vidéos
suit une structure à trois niveaux (partition, sujet, clip) et mélange deux
formats de conteneur. Plutôt que de reconstruire un chemin à chaque appel, le
module de résolution construit \emph{une seule fois} un index en mémoire associant
chaque nom de fichier à son chemin absolu, puis tolère l'absence d'extension dans
l'identifiant en essayant successivement les deux formats connus. Cette
indexation initiale rend le coût de résolution négligeable et, surtout, rend
explicite le cas d'un clip annoté dont la vidéo est absente du disque : la
situation est retournée comme un état d'erreur nommé plutôt que comme une
exception.

\paragraph{Extraction des régions de visage.} La perception faciale ne travaille
pas sur l'image entière mais sur la région du visage, isolée par
MediaPipe Face Detector (BlazeFace). L'implémentation enchaîne trois
tentatives --- courte portée, puis longue portée, puis image miroir --- et
retient la boîte de plus grande aire, la méthode effectivement employée
étant conservée dans la trace. Ce dispositif répond à une caractéristique
du corpus : les apprenants filmés en situation d'ennui ou de distraction
détournent fréquemment la tête, et une détection frontale seule échouerait
précisément sur les cas pédagogiquement les plus intéressants. Le
tableau~\ref{tab:params-extraction} rassemble les paramètres retenus.

\begin{table}[H]
\centering
\small
\renewcommand{\arraystretch}{1.2}
\begin{tabular}{|>{\raggedright\arraybackslash}p{4.6cm}|>{\raggedright\arraybackslash}p{2.5cm}|>{\raggedright\arraybackslash}p{5.8cm}|}
\hline
\rowcolor{gray!20}
\textbf{Opération} & \textbf{Valeur} & \textbf{Motif du réglage} \\
\hline
Images analysées par clip & 5 &
Compromis entre robustesse et coût sur processeur
(section~\ref{subsec:stack}). \\
\hline
Position des images & régulière, hors bords &
Évite les premières et dernières images, souvent floues ou en transition. \\
\hline
Égalisation d'histogramme & non (MediaPipe) &
BlazeFace opère sur l'image RGB ; l'égalisation Haar n'est plus requise. \\
\hline
Modèle court / long & BlazeFace &
Court : webcam proche (DAiSEE) ; long : plan plus large (IEMOCAP). \\
\hline
Confiance minimale & 0{,}5 &
Filtre les détections trop incertaines. \\
\hline
Taille minimale de visage & 40 $\times$ 40 pixels &
Écarte les régions trop petites pour porter une expression exploitable. \\
\hline
Marge autour de la boîte & 30 \% &
Conserve le front et le menton, informatifs pour l'ennui et la confusion. \\
\hline
\end{tabular}
\caption{Paramètres de la chaîne d'extraction visuelle appliquée aux clips
\textsc{DAiSEE}.}
\label{tab:params-extraction}
\end{table}

\paragraph{Nettoyage et politique d'exclusion.} Le nettoyage retenu est
\emph{déclaratif} plutôt que silencieux. Un clip est écarté du jeu d'entraînement
dans quatre cas seulement : vidéo introuvable, vidéo illisible, vidéo sans image
décodable, ou échec du classifieur d'expressions ; chacun de ces cas porte un code
distinct. Aucune suppression de doublons ni filtrage sur la durée n'a été appliqué,
les clips de \textsc{DAiSEE} étant tous calibrés à dix secondes. Aucun
rééquilibrage artificiel des classes n'a été effectué non plus : le déséquilibre
massif du corpus, déjà documenté au chapitre~\ref{chap:conception}
(figure~\ref{fig:distribution}), est traité en aval par une pondération inverse des
classes lors de l'apprentissage et par le choix d'une métrique moyennée par classe,
ce qui préserve la distribution réelle du phénomène observé.

\paragraph{Mise en cache des traits.} L'extraction des traits faciaux est
l'opération la plus coûteuse de la chaîne : elle suppose de décoder la vidéo, de
détecter le visage sur cinq images et d'appeler le réseau de classification. Pour
que l'apprentissage du mapping puis les analyses de sensibilité puissent être
relancés sans repayer ce coût, les vecteurs de sept probabilités sont écrits une
fois pour toutes dans un fichier de cache accompagnant chaque clip de son état
annoté. Le cache constitué compte 1\,200 clips ; il reproduit la distribution
déséquilibrée du corpus (885 clips d'engagement contre 42 de frustration), ce qui
en fait aussi un artefact de documentation du problème d'apprentissage traité.

\subsection{Extraction et alignement du corpus IEMOCAP}
\label{subsec:impl-iemocap}

\paragraph{Analyse des fichiers d'annotation.} \textsc{IEMOCAP} distribue ses
annotations sous forme de fichiers textuels par dialogue, où chaque énoncé occupe
une ligne mêlant bornes temporelles, identifiant, étiquette émotionnelle et
coordonnées d'activation, de valence et de dominance. Une expression régulière
unique extrait ces sept champs et alimente une table homogène. Le résultat,
matérialisé dans un fichier unique de \textbf{10\,039 énoncés} couvrant les cinq
sessions, remplace définitivement la lecture des fichiers d'origine : tous les
traitements ultérieurs — projection vers la taxonomie d'engagement, analyse de
sensibilité, entraînement de la tête vocale — partent de cette table. La
projection des étiquettes vers les quatre états \textsc{DAiSEE} suit la
correspondance de référence définie au chapitre~\ref{chap:conception}
(tableau~\ref{tab:mapping-iemocap}) ; elle est implémentée comme un dictionnaire
isolé, ce qui a permis l'analyse de sensibilité présentée en
section~\ref{subsec:sensibilite-mapping} par simple substitution d'une entrée.

\paragraph{Localisation du signal vocal.} Le corpus ne fournit pas un fichier
audio par énoncé mais un fichier par dialogue. Le chemin du dialogue est reconstruit
à partir de l'identifiant de l'énoncé, dont le suffixe de locuteur est retiré par
expression régulière, puis le segment annoté est chargé à la position et pour la
durée indiquées, en monophonie et à 16~kHz — la fréquence attendue par l'encodeur
vocal pré-entraîné. Une étape d'enrichissement a ajouté à la table le chemin
résolu et un indicateur d'existence pour chaque énoncé, ce qui permet de vérifier
en une lecture que la couverture audio est complète avant tout entraînement.

\paragraph{Appariement voix--scénario.} \textsc{DAiSEE} et \textsc{IEMOCAP} ne
partageant aucun sujet, l'appariement ne peut être qu'une simulation, dont le
principe a été posé au chapitre~\ref{chap:conception}
(section~\ref{subsec:appariement}). Son implémentation est volontairement
minimale et déterministe : les énoncés dont l'état projeté correspond à l'état
dominant du scénario sont regroupés, puis l'énoncé retenu est désigné par le reste
de la division de l'identifiant du scénario par la taille du groupe. Le même
scénario reçoit donc toujours le même énoncé, propriété indispensable pour que
deux exécutions du banc d'essai soient comparables, et l'appariement est signalé
comme tel dans la réponse de l'API afin qu'il ne puisse jamais être confondu avec
une captation réellement synchrone.

\paragraph{Sous-corpus audiovisuel réellement aligné.} Pour disposer malgré tout
d'une mesure où visage et voix proviennent de la même personne au même instant,
un second artefact a été construit à l'intérieur d'\textsc{IEMOCAP} lui-même
(section~\ref{subsec:option-a}). Seule la première session dispose des
enregistrements vidéo ; ceux-ci filment deux locuteurs côte à côte dans une image
de 720 $\times$ 480 pixels, ce qui impose de déterminer quel côté correspond au
locuteur de l'énoncé. Le genre du locuteur, lisible dans l'identifiant, est
comparé au genre annoncé dans le nom du dialogue, et la moitié d'image
correspondante est découpée avant application de la même chaîne de détection de
visage que pour \textsc{DAiSEE}. Après filtrage des énoncés trop courts et des
étiquettes indéterminées, un échantillon équilibré de vingt énoncés par état
constitue un sous-corpus de \textbf{80 paires audiovisuelles} sur lequel les deux
modalités et leurs fusions seront comparées au chapitre~\ref{chap:validation}.

\paragraph{Constitution du jeu d'entraînement vocal.} La tête de classification
vocale est apprise sur ce même référentiel, filtré selon trois critères explicites :
exclusion des étiquettes indéterminées, durée minimale d'une seconde et présence
effective du fichier audio. Un tirage par classe ramène ensuite les quatre états à
des effectifs voisins, ce qui compense la rareté de la confusion dans le corpus.
Les segments sont recadrés ou complétés à une longueur fixe de quatre secondes,
l'encodeur est gelé et seul le classifieur final est optimisé, sur deux époques et
avec un partage stratifié à 80 / 20. Tous les tirages dépendent d'une graine unique,
consignée dans le fichier de métriques produit à côté du modèle.

\subsection{Création de scénarios d'apprentissage}
\label{subsec:impl-scenarios}

\paragraph{Du clip annoté au scénario pédagogique.} Un clip de \textsc{DAiSEE} est
un objet perceptif : dix secondes de vidéo et quatre scores. Un \emph{scénario
d'apprentissage} est un objet pédagogique : un état d'engagement situé dans un
contexte de cours, auquel le système doit répondre. Le passage de l'un à l'autre
est opéré par deux fonctions volontairement séparées. La première détermine l'état
dominant du clip, en levant les égalités par un ordre de priorité pédagogique qui
place la frustration avant la confusion, puis l'ennui, puis l'engagement : à scores
égaux, c'est le besoin le plus urgent qui est retenu, jamais l'état le plus
favorable. La seconde attribue au clip un \emph{patron}, c'est-à-dire une des six
configurations pédagogiques définies au chapitre~\ref{chap:conception}
(tableau~\ref{tab:patterns}), en traitant d'abord les deux cas composites — la
situation critique où au moins trois dimensions dépassent le seuil, et
l'engagement faible — avant de retomber sur le patron associé à l'état dominant.

\paragraph{Sélection des cinquante scénarios.} Un échantillon de cinquante clips a
été constitué à partir du seul sous-ensemble de conception, par tirages
indépendants et reproductibles pour chaque patron
(tableau~\ref{tab:composition-scenarios}). Les effectifs ne sont pas
proportionnels à la fréquence des patrons dans le corpus : ils sont
\emph{délibérément déséquilibrés en faveur des états problématiques}. Un tirage
proportionnel aurait produit un échantillon composé presque exclusivement de clips
engagés, sur lesquels le système n'a par construction rien à décider ; le
sur-échantillonnage des situations de frustration, de confusion et d'état critique
est ce qui rend l'échantillon informatif pour évaluer la chaîne de décision. Ce
choix est aussi une limite d'interprétation, rappelée au
chapitre~\ref{chap:validation} : les taux mesurés sur ces cinquante scénarios ne
sont pas des taux attendus en population.

\begin{table}[H]
\centering
\small
\renewcommand{\arraystretch}{1.2}
\begin{tabular}{|>{\raggedright\arraybackslash}p{3.4cm}|>{\raggedright\arraybackslash}p{5.4cm}|c|}
\hline
\rowcolor{gray!20}
\textbf{Patron} & \textbf{Critère de tirage} & \textbf{Effectif} \\
\hline
Stable & engagement $\geq 2$ & 15 \\
\hline
Ennui élevé & ennui $\geq 2$ & 10 \\
\hline
Frustration & frustration $\geq 2$ & 8 \\
\hline
Confusion & confusion $\geq 2$ & 7 \\
\hline
Engagement faible & engagement $\leq 1$ & 6 \\
\hline
État critique & confusion $\geq 2$ \emph{et} frustration $\geq 2$ & 4 \\
\hline
\textbf{Total} & & \textbf{50} \\
\hline
\end{tabular}
\caption{Composition du référentiel de scénarios : critère de tirage et effectif
par patron pédagogique.}
\label{tab:composition-scenarios}
\end{table}

\paragraph{Un référentiel recalculé plutôt que figé.} Le fichier de scénarios
conserve, outre l'identifiant du clip et ses quatre scores, le patron obtenu au
moment du tirage. Ce patron est néanmoins \textbf{recalculé au chargement} par le
serveur à partir des scores. Cette redondance apparente est un garde-fou : elle
garantit qu'une modification ultérieure des règles de catégorisation se propage
immédiatement au référentiel, au lieu de laisser cohabiter des étiquettes
héritées d'une version antérieure du code. Un jeu de tests dédié vérifie d'ailleurs
que, pour chacune des cinquante lignes, le patron stocké reste cohérent avec l'état
dominant. Un second fichier dérivé associe à chaque scénario la stratégie et le
type de contenu retenus, servant de table de référence lisible pour la discussion
des résultats.

\subsection{Segmentation des sessions}
\label{subsec:impl-segmentation}

\paragraph{Trois granularités emboîtées.} La simulation manipule trois niveaux de
découpage temporel. Le \textbf{clip} de dix secondes est l'unité de perception et
de décision : c'est le grain auquel l'API répond et auquel une adaptation est
éventuellement produite. La \textbf{fenêtre glissante} est l'unité d'analyse
dynamique : les règles de déclenchement examinent les deux ou quatre derniers clips
selon la règle considérée (section~\ref{subsec:impl-transitions}). La
\textbf{séquence} est l'unité de parcours : dix séquences de quatre à sept clips
ont été assemblées hors ligne à partir du référentiel, chacune portant un titre et
un patron dominant — engagement stable, ennui progressif, frustration montante,
confusion avec blocage, état critique multi-états. Ces séquences ne sont pas
tirées au hasard : elles sont construites pour exposer des \emph{trajectoires}
plutôt que des états isolés, seule manière d'observer le comportement des règles
temporelles.

\paragraph{Persistance des sessions.} Chaque adaptation effectivement produite est
enregistrée dans la base de sessions avec l'identifiant du scénario, l'état
détecté, son intensité, la stratégie, le type de contenu, le texte généré et un
horodatage. Le choix d'une table plate, sans identifiant de parcours, mérite d'être
assumé explicitement : la reconstitution d'une trajectoire s'effectue par lecture
ordonnée des enregistrements, ce qui suffit à un prototype mono-utilisateur mais ne
permettrait pas de distinguer plusieurs apprenants simultanés. Le journal de
décision, lui, conserve pour chaque clip analysé les éléments nécessaires à
l'analyse scientifique — état fusionné, type de conflit entre modalités, patron
retenu, indicateurs de génération — sous une forme en ajout seul qui interdit la
réécriture rétrospective d'une décision.

%==============================================================================
\section{Implémentation du modèle d'adaptation}
\label{sec:impl-modele}

Le modèle d'adaptation transforme un état perçu en une intention pédagogique. Sa
formulation conceptuelle occupe les sections~\ref{sec:perception}
et~\ref{sec:modele} du chapitre~\ref{chap:conception} ; il s'agit ici d'exposer
sa traduction en code,
c'est-à-dire l'ordre exact d'évaluation des règles, les seuils effectivement
programmés et le traitement des cas où deux règles pourraient s'appliquer. Trois
composants sont concernés : la fusion des modalités, la sélection de la stratégie
et le patron de prompt, et la prise en compte de la dynamique temporelle.

\subsection{Règles de décision}
\label{subsec:impl-regles}

\paragraph{Fusion des deux modalités.} La fusion est implémentée en deux temps
séparés : \emph{qualifier} la relation entre les deux perceptions, puis
\emph{résoudre} le désaccord éventuel. La qualification distingue six situations,
selon que les deux modalités désignent le même état ou non. Lorsqu'elles
s'accordent, l'écart d'intensité est comparé à un seuil de tolérance d'un niveau :
en deçà, l'accord est jugé suffisant et l'intensité finale est la moyenne pondérée
des deux, avec un poids de 0{,}60 pour le visage et de 0{,}40 pour la voix ; au-delà,
un conflit d'intensité est déclaré. Lorsqu'elles désignent des états différents, la
sévérité du conflit est mesurée par l'écart entre les priorités pédagogiques des
deux états, ce qui distingue un désaccord mineur — par exemple entre ennui et
confusion — d'un désaccord majeur entre frustration et engagement. La résolution
applique alors la stratégie sélectionnée par configuration
(section~\ref{subsec:choix-archi}) : la priorité pédagogique retient l'état le plus
urgent, la stratégie de confiance retient celui dont le score est le plus élevé.
Un point important pour l'interprétation des résultats : quelle que soit la
stratégie, la sortie conserve le type de conflit et la règle appliquée, si bien
qu'une décision peut toujours être rejouée et expliquée.

\paragraph{Deux familles de fusion, deux usages.} La distinction établie au
chapitre~\ref{chap:conception} entre fusion par résolution de conflit et fusion par
vote pondéré (section~\ref{subsec:vote}) se traduit par deux fonctions
coexistantes. La première, employée par le serveur, décide d'abord de l'état puis
pondère l'intensité : le poids n'influence donc pas le choix de l'état. La seconde
combine les distributions de probabilité des quatre états avant d'en prendre
l'argument maximal, de sorte que le poids agit cette fois sur l'état retenu ; elle
est mobilisée par le banc d'essai sur le sous-corpus audiovisuel aligné, où elle
permet le balayage des poids de modalité rapporté au chapitre~\ref{chap:validation}.
Faire cohabiter les deux implémentations, plutôt que d'en généraliser une, était la
seule manière de comparer les deux familles sur les mêmes données sans modifier le
comportement du prototype.

\paragraph{Cascade de sélection de la stratégie.} La décision proprement
pédagogique est une cascade de conditions évaluées dans un ordre fixe, la première
satisfaite l'emportant (tableau~\ref{tab:cascade-decision}). Cet ordre \emph{est}
la règle : il encode une hiérarchie de besoins où le blocage combinant confusion et
frustration passe avant le désengagement total, lui-même avant la frustration
seule, la confusion, l'ennui, puis l'engagement. Écrire la même logique comme une
suite de conditions indépendantes aurait laissé indéterminé le cas — fréquent dans
le corpus — d'un clip présentant simultanément plusieurs scores élevés.

\begin{table}[H]
\centering
\small
\renewcommand{\arraystretch}{1.2}
\begin{tabular}{|c|>{\raggedright\arraybackslash}p{3.5cm}|>{\raggedright\arraybackslash}p{5.0cm}|c|}
\hline
\rowcolor{gray!20}
\textbf{Rang} & \textbf{Condition} & \textbf{Stratégie retenue} &
\textbf{Type} \\
\hline
1 & confusion $\geq 3$ et frustration $\geq 2$ & Aide de déblocage immédiate &
A+D \\
\hline
2 & engagement $= 0$ & Réengagement urgent & B \\
\hline
3 & frustration $\geq 2$ & Encouragement et simplification & D+A \\
\hline
4 & confusion $\geq 2$ & Explication alternative & A \\
\hline
5 & ennui $\geq 2$ & Exercice avancé & C \\
\hline
6 & engagement $\leq 1$ & Exemple concret stimulant & B \\
\hline
7 & ennui $= 1$ & Application pratique & B \\
\hline
8 & engagement $\geq 2$ & Renforcement positif & D \\
\hline
9 & aucun des cas précédents & Continuation normale & aucune \\
\hline
\end{tabular}
\caption{Cascade de décision telle qu'implémentée : la première condition
satisfaite détermine la stratégie et le type de contenu.}
\label{tab:cascade-decision}
\end{table}

\subsection{Algorithme de sélection de stratégies}
\label{subsec:impl-selection}

\paragraph{Une sélection en deux étages.} Le type de contenu issu de la cascade ne
suffit pas à lancer une génération : il désigne une famille d'intervention, non un
patron de prompt. Un second étage traduit le couple formé par ce type et le
contexte du clip en l'un des vingt patrons de la bibliothèque
(tableau~\ref{tab:templates}). Ce découplage a une raison précise : il permet de
faire évoluer la bibliothèque de patrons — en ajouter, en reformuler — sans
toucher aux règles pédagogiques, et réciproquement de modifier un seuil de décision
sans réécrire de prompt. L'algorithme complet se déroule ainsi :

\begin{enumerate}[itemsep=2pt, topsep=3pt]
\item les quatre scores sont lus dans l'état fusionné ;
\item la cascade du tableau~\ref{tab:cascade-decision} désigne un type de contenu
      et une stratégie ;
\item si la dynamique temporelle conclut à l'absence de besoin d'intervention, la
      décision est court-circuitée
      (section~\ref{subsec:impl-transitions}) ;
\item sinon, le type est raffiné en un patron précis à l'aide de l'intensité de
      l'état, du patron du clip et des scores individuels ;
\item les métadonnées du patron — objectif, longueur cible, format, ton — sont
      lues dans la bibliothèque et jointes à la décision ;
\item l'objet de décision, qui porte à la fois le résultat et ses déterminants,
      est transmis au module de génération.
\end{enumerate}

\paragraph{Règles de raffinement.} Le raffinement obéit à des règles courtes et
lisibles. Le type composite associant explication et soutien conduit directement au
patron d'aide au déblocage. Le soutien guidé se scinde selon l'intensité de la
frustration, un niveau maximal appelant un patron de soutien plus enveloppant
qu'un patron d'encouragement standard. L'explication alternative se décline en
trois variantes : une micro-explication décomposée lorsque la confusion est
maximale, une explication reformulée lorsque le patron du clip confirme la
confusion, une explication narrative dans les autres cas. L'exemple concret
distingue le réengagement d'un apprenant totalement décroché de la simple
stimulation d'un apprenant peu engagé. Enfin, un cas par défaut garantit qu'aucune
combinaison de scores ne peut laisser la sélection sans patron.

\paragraph{Ce que produit la cascade sur le référentiel.} Appliquée aux cinquante
scénarios, la chaîne complète ne mobilise \textbf{que onze des vingt patrons}, la
distribution étant fortement concentrée sur l'encouragement en cas de frustration
modérée et sur l'exercice avancé en cas d'ennui élevé. Ce résultat, obtenu et
consigné par le banc d'essai, n'est pas un défaut d'implémentation mais la
conséquence conjointe de la composition du référentiel et de l'ordre de la
cascade : les patrons dédiés aux situations rares ne peuvent être atteints que par
des configurations de scores que l'échantillon contient peu ou pas. Il constitue à
ce titre une observation de premier plan pour la discussion du
chapitre~\ref{chap:validation}, et une justification du corpus complémentaire
généré en section~\ref{sec:impl-corpus}, qui couvre dix-neuf patrons.

\subsection{Gestion des transitions d'état}
\label{subsec:impl-transitions}

\paragraph{Quatre degrés d'urgence.} La dynamique temporelle est implémentée dans
un module distinct, qui reçoit l'historique récent et retourne quatre éléments :
faut-il intervenir, avec quel degré d'urgence, en vertu de quelle règle nommée, et
depuis combien de clips l'état persiste. Les degrés — immédiate, rapide, différée,
aucune — sont évalués dans cet ordre, si bien qu'une situation de blocage aigu ne
peut jamais être ramenée à une intervention différée par une règle de rang
inférieur. La détection immédiate ne considère que le clip courant, dans deux cas
seulement : confusion maximale accompagnée de frustration, ou désengagement total.
L'urgence rapide exige la persistance d'un seuil sur \textbf{deux clips
consécutifs}, ce qui filtre les fluctuations d'une seule fenêtre de dix secondes.
L'urgence différée observe \textbf{quatre clips}, pour capter les états faibles mais
durables — engagement bas prolongé, ennui de faible intensité qui s'installe — que
la règle rapide laisserait passer.

\paragraph{Le cas de la non-intervention.} La règle la plus importante pour la
crédibilité pédagogique du système est celle qui conclut à \emph{ne rien faire} :
lorsque l'engagement se maintient à un niveau satisfaisant sur les deux derniers
clips, aucune adaptation n'est déclenchée. Son implémentation traverse toute la
chaîne. Le moteur de décision force alors le type de contenu à « aucune », remplace
la stratégie par une continuation normale et marque explicitement la décision comme
non déclenchée ; l'orchestrateur n'appelle pas le service de génération et
substitue au contenu un objet portant un indicateur de saut ; enfin, ni le journal
de décision ni la base de sessions ne reçoivent d'enregistrement. Une
non-intervention ne consomme donc aucun appel au modèle de langue et reste
distinguable, dans les traces, d'une intervention qui aurait échoué.

\paragraph{Ce qui est implémenté, ce qui ne l'est pas.} Deux limites doivent être
énoncées ici plutôt que découvertes au chapitre~\ref{chap:validation}. D'une part,
le module temporel est complètement implémenté et couvert par des tests unitaires
qui vérifient chacun de ses cas — déclenchement immédiat sur blocage,
déclenchement immédiat sur désengagement, non-intervention sur engagement stable,
déclenchement rapide sur ennui consécutif, court-circuitage de la décision — mais
le point d'entrée exposé par l'API analyse un clip isolé : la dynamique n'est
mobilisée que lorsqu'une trajectoire est rejouée à partir des séquences
reconstituées, et non lors d'une analyse ponctuelle depuis l'interface. D'autre
part, le moteur de décision est \textbf{sans mémoire d'une décision à la suivante} :
il n'implémente ni hystérésis, ni délai de garde entre deux interventions, ni
mécanisme évitant la répétition d'un même patron sur des clips voisins. Sur un
parcours réel, ces absences produiraient un risque d'oscillation et de saturation
de l'apprenant ; elles constituent la première piste d'amélioration formulée au
chapitre~\ref{chap:conclusion}.

%==============================================================================
\section{Implémentation du module de génération adaptative}
\label{sec:impl-generation}

La décision étant prise, il reste à produire le contenu. Cette section décrit la
configuration du service de génération, la bibliothèque de patrons qui structure
les demandes, l'assemblage du prompt effectivement envoyé, et les dispositifs de
contrôle appliqués aux textes obtenus.

\subsection{Configuration du LLM}
\label{subsec:impl-llm}

\paragraph{Choix du service et du modèle.} La génération est confiée à un modèle
\textsc{LLaMA}~3.1 de huit milliards de paramètres \cite{touvron2023llama},
interrogé à distance via l'interface de programmation de Groq. Ce choix résulte des
contraintes énoncées en section~\ref{sec:infrastructure} : un modèle propriétaire
de très grande taille aurait supposé un budget d'appels incompatible avec la
génération d'un corpus de plusieurs centaines de contenus, et un modèle hébergé
localement était exclu par l'absence d'accélérateur. Le modèle est utilisé
\textbf{sans réglage fin} : toute l'adaptation passe par le prompt, ce qui rend
chaque contenu explicable par la conjonction d'un état détecté, d'un patron et d'un
contexte documentaire, et rend le dispositif transposable à un autre fournisseur
sans réapprentissage.

\paragraph{Paramètres de décodage.} Les paramètres retenus
(tableau~\ref{tab:params-llm}) traduisent une exigence de sobriété plutôt que de
créativité. Une température basse limite la variabilité entre deux générations
comparables, ce qui est indispensable pour que le corpus produit soit
interprétable. Un budget de sortie borné empêche les productions
interminables et maintient les contenus dans l'ordre de grandeur attendu par les
patrons. Le message de rôle système, enfin, formule trois interdits : produire en
français, n'utiliser que les informations du contexte de référence fourni, et
déclarer explicitement l'absence d'information plutôt que de l'inventer — première
ligne de défense contre les hallucinations, complétée par le module RAG
(ancrage documentaire) décrit en section~\ref{subsec:impl-pipeline-gen}.

\begin{table}[H]
\centering
\small
\renewcommand{\arraystretch}{1.2}
\begin{tabular}{|>{\raggedright\arraybackslash}p{4.0cm}|>{\raggedright\arraybackslash}p{2.7cm}|>{\raggedright\arraybackslash}p{6.2cm}|}
\hline
\rowcolor{gray!20}
\textbf{Paramètre} & \textbf{Valeur} & \textbf{Effet recherché} \\
\hline
Modèle & \texttt{openai/gpt-oss-20b} &
Remplaçant Groq de \texttt{llama-3.1-8b-instant} (retiré le 16/08/2026) :
coût faible et latence faible, au prix d'une capacité inférieure aux très grands
modèles. \\
\hline
Température (contenu) & 0{,}3 &
Contenus reproductibles à contexte constant. \\
\hline
Température (tuteur) & 0{,}25 à 0{,}55 &
Modulée selon l'état : basse en confusion, plus haute en situation d'ennui. \\
\hline
Budget de sortie & 600 unités &
Borne la longueur des contenus adaptatifs. \\
\hline
Budget de sortie (tuteur) & 400 à 650 unités &
Ajusté selon la demande de l'apprenant (exemple, exercice, détail). \\
\hline
Tentatives & 4, espacées de 2, 5 puis 8~s &
Absorbe les limitations de débit du service gratuit. \\
\hline
Fenêtre de dialogue & 10 derniers messages &
Borne le contexte conversationnel et le coût par échange. \\
\hline
Clé d'accès & variable d'environnement &
Aucun secret dans le dépôt. \\
\hline
\end{tabular}
\caption{Configuration du service de génération et du tuteur conversationnel.}
\label{tab:params-llm}
\end{table}

\paragraph{Comportement en cas d'indisponibilité.} La gestion des défaillances est
partielle et il importe de le dire. Les dépassements de quota sont explicitement
rattrapés par des tentatives espacées, mécanisme qui a rendu possible la génération
du corpus décrite en section~\ref{sec:impl-corpus}. En revanche, aucun délai
d'expiration n'est fixé sur l'appel distant et les autres erreurs de service ne
sont pas retentées : elles remontent jusqu'au point d'entrée de l'API, qui répond
par un message d'erreur explicite. Le prototype ne propose donc pas de contenu de
repli hors ligne — un choix assumé, puisqu'un contenu produit sans le modèle
n'aurait pas la même nature et fausserait l'analyse du corpus.

\subsection{Bibliothèque de 20 templates de prompts}
\label{subsec:impl-templates}

\paragraph{Séparation des métadonnées et des instructions.} La bibliothèque des
vingt patrons introduite au chapitre~\ref{chap:conception}
(tableau~\ref{tab:templates}) est implémentée en deux structures distinctes. La
première est une table de \textbf{métadonnées} : pour chaque patron, sa famille,
sa longueur cible en mots, son format attendu, son ton et son objectif pédagogique.
Ces attributs ne servent pas seulement à documenter : ils sont injectés dans la
section des contraintes du prompt, exposés par l'API et réutilisés par l'audit
automatique comme référence de longueur (section~\ref{subsec:impl-postproc}). La
seconde est une table de \textbf{constructeurs d'instructions}, une fonction par
patron, qui produit le corps de la demande en y insérant les variables de contexte :
concept visé, intensité de l'état, intensité spécifique de la frustration, durée de
persistance et module en cours. Un patron inconnu retombe sur un patron de
référence, de sorte qu'une faute de frappe dans un identifiant ne provoque jamais
d'interruption.

\paragraph{Anatomie d'un patron.} Les instructions ne demandent pas un texte libre
mais une structure. Le patron d'aide au déblocage, mobilisé lorsque confusion et
frustration se combinent, demande ainsi quatre éléments dans un ordre imposé : une
phrase rassurante qui normalise la difficulté, l'identification du point de blocage
probable, une approche radicalement différente de l'explication initiale, puis un
micro-exercice guidé très simple — le tout sur environ deux cents mots et sur un ton
déculpabilisant. Le patron de félicitation, à l'autre extrémité de la bibliothèque,
demande de célébrer la réussite, d'expliciter la compétence acquise, de faire le
lien avec la suite du parcours et de renforcer la confiance, en cent cinquante mots.
Cette forme numérotée n'est pas un détail stylistique : elle rend le contenu
généré \emph{vérifiable}, puisqu'un auditeur — automatique ou humain — peut
constater la présence ou l'absence de chaque élément attendu.

\subsection{Pipeline de génération}
\label{subsec:impl-pipeline-gen}

\paragraph{Ancrage documentaire (RAG).} Avant toute génération, un contexte de
référence est constitué par \textbf{récupération augmentée}
(\emph{Retrieval-Augmented Generation}, RAG) \cite{lewis2020retrieval}. Neuf fiches
documentaires — définitions, exemples, exercices sur les algorithmes de tri, les
réseaux de neurones et les structures de données, auxquelles s'ajoutent deux fiches
sur l'encouragement et les méthodes d'apprentissage — sont encodées par un modèle
de phrases multilingue et indexées dans une base vectorielle persistante utilisant
la similarité cosinus. Les documents étant courts, aucun découpage en fragments n'a
été nécessaire. La requête n'est pas le concept seul : elle est enrichie de l'état
détecté, de sorte qu'un apprenant frustré et un apprenant curieux sur le même
concept ne reçoivent pas nécessairement le même appui documentaire. Les trois
documents les plus proches sont retenus, à condition de dépasser un seuil de
similarité ; en deçà, ils sont écartés et un contexte générique est fourni, ce qui
évite d'ancrer un contenu sur une source hors sujet.

\paragraph{Assemblage du prompt.} Le prompt transmis au modèle est une composition
de cinq blocs ordonnés (figure~\ref{fig:prompt}). Cet ordre suit une progression
délibérée : la référence factuelle d'abord, la situation pédagogique ensuite,
l'état affectif et sa durée, puis seulement la tâche et ses contraintes. La durée
de persistance de l'état est convertie en secondes dans le prompt, afin que le
modèle raisonne sur une réalité temporelle et non sur un compteur abstrait de
clips.

\begin{figure}[H]
\centering
\begin{tikzpicture}[
  font=\footnotesize,
  blk/.style={rectangle, draw=black, thick, align=left, text width=11.4cm,
              minimum height=0.8cm, inner sep=4pt, fill=white},
  arr/.style={-{Stealth[length=2.4mm]}, thick}
]
\node[blk] (b1) {\textbf{1. Contexte de référence} — trois documents récupérés
  (type, source, contenu) ou contexte générique};
\node[blk, below=0.28cm of b1] (b2) {\textbf{2. Contexte pédagogique} — sujet,
  module, concept, niveau de l'apprenant, objectif};
\node[blk, below=0.28cm of b2] (b3) {\textbf{3. État d'engagement} — état,
  intensité sur 3, durée de persistance, patron, moment du cours};
\node[blk, below=0.28cm of b3] (b4) {\textbf{4. Tâche et contraintes} — famille de
  contenu, longueur cible, format, ton};
\node[blk, below=0.28cm of b4] (b5) {\textbf{5. Instructions du patron} — structure
  numérotée propre au patron retenu};
\node[rectangle, rounded corners, draw=black, thick, fill=black!8, below=0.45cm of b5,
      text width=11.4cm, align=center, minimum height=0.8cm] (out)
  {Contenu adaptatif en français + indicateurs (nombre de mots, sources
   documentaires)};
\draw[arr] (b5) -- (out);
\end{tikzpicture}
\caption{Composition du prompt de génération : cinq blocs assemblés dans un ordre
fixe, précédés d'un message de rôle système interdisant l'invention de faits.}
\label{fig:prompt}
\end{figure}

\paragraph{Sortie et traçabilité.} La réponse retournée par le module de génération
ne contient pas seulement le texte. Elle porte l'état et l'intensité qui l'ont
motivée, le patron employé, l'objectif visé, le nombre de mots produits, un
indicateur d'utilisation effective de la récupération documentaire, le nombre et les
sources des documents injectés, et le cas échéant le message d'erreur de la
récupération. L'interface expose ces éléments à l'utilisateur et le journal en
conserve les principaux : un contenu du corpus peut ainsi être rattaché, après
coup, à la décision et aux sources qui l'ont produit.

\paragraph{Tuteur conversationnel.} Le même dispositif alimente un second mode
d'interaction. Le tuteur reçoit l'état détecté et adapte trois éléments : le ton et
la consigne pédagogique, propres à chaque état ; la température de décodage, plus
basse en situation de confusion pour privilégier la précision, plus haute en
situation d'ennui pour autoriser la variété ; et le budget de réponse, étendu
lorsque la question de l'apprenant appelle un exemple, un exercice ou un
approfondissement. Le contenu adaptatif initial est réinjecté comme premier tour de
parole, de sorte que le dialogue prolonge l'adaptation au lieu de repartir de zéro.

\subsection{Post-processing et validation}
\label{subsec:impl-postproc}

\paragraph{Un contrôle a posteriori, hors de la boucle de génération.} Le module de
génération n'applique aucune transformation au texte reçu, hormis le comptage des
mots : ni troncature, ni réécriture, ni nouvelle tentative en cas de sortie jugée
insatisfaisante. Ce choix est méthodologique. Un post-traitement corrigeant
silencieusement les productions rendrait impossible la mesure de la qualité
intrinsèque du dispositif : on n'évaluerait plus la génération adaptative, mais la
génération augmentée d'un correcteur. Le contrôle est donc reporté sur un
\textbf{audit hors ligne} appliqué à l'ensemble du corpus, dont les verdicts
alimentent une régénération ciblée (section~\ref{sec:impl-corpus}).

\paragraph{Critères de l'audit automatique.} L'audit vérifie trois familles de
propriétés (tableau~\ref{tab:audit-criteres}) : la conformité de longueur, évaluée
par rapport à une plage tolérée autour de la longueur cible du patron ; la
cohérence entre l'émotion enregistrée et la famille du patron ; et la présence de
marqueurs lexicaux attendus, détectés par expressions régulières — marqueurs
d'explication pour la famille A, de consigne d'exercice pour la famille C, de
soutien pour les patrons de frustration. Chaque anomalie détectée porte un code
nommé et retire vingt-cinq points à un score de qualité initialisé à cent, un
contenu sans anomalie étant déclaré conforme. Cette granularité rend l'audit
actionnable : le rapport indique non pas qu'un contenu est mauvais, mais quel
critère précis il enfreint.

\begin{table}[H]
\centering
\small
\renewcommand{\arraystretch}{1.2}
\begin{tabular}{|>{\raggedright\arraybackslash}p{5.3cm}|>{\raggedright\arraybackslash}p{7.6cm}|}
\hline
\rowcolor{gray!20}
\textbf{Code d'anomalie} & \textbf{Propriété vérifiée} \\
\hline
\texttt{trop\_court} · \texttt{trop\_long} &
Longueur du contenu au regard de la cible déclarée par le patron. \\
\hline
\texttt{emotion\_\allowbreak template\_suspect} &
Compatibilité entre l'émotion enregistrée et la famille du patron. \\
\hline
\texttt{manque\_marqueurs\_\allowbreak explication} &
Présence de marqueurs de reformulation dans un contenu de famille A. \\
\hline
\texttt{manque\_marqueurs\_\allowbreak exercice} &
Présence d'une consigne d'exercice dans un contenu de famille C. \\
\hline
\texttt{ton\_trop\_positif\_\allowbreak pour\_frustration} &
Absence de registre de félicitation dans un contenu de soutien. \\
\hline
\texttt{confusion\_sans\_ton\_\allowbreak rassurant} &
Présence d'un registre rassurant dans un contenu de déblocage. \\
\hline
\texttt{template\_inconnu} &
Existence du patron déclaré dans la bibliothèque. \\
\hline
\end{tabular}
\caption{Critères de l'audit automatique du contenu généré et codes d'anomalie
associés.}
\label{tab:audit-criteres}
\end{table}

\paragraph{Indicateurs de lisibilité.} Un second script calcule, indépendamment de
l'audit de cohérence, des indicateurs de forme : nombre de mots et de phrases,
longueur moyenne des phrases, indice de lisibilité approché par comptage
syllabique, et taux de répétition défini comme le complément du rapport entre mots
distincts et mots totaux. Ces indicateurs ne portent aucun jugement pédagogique ; ils
fournissent une caractérisation objective du corpus, exploitée au
chapitre~\ref{chap:validation} en complément de l'appréciation humaine, dont ils ne
peuvent tenir lieu.

%==============================================================================
\section{Développement du prototype de simulation}
\label{sec:impl-prototype}

Le prototype n'a pas pour fonction de démontrer une faisabilité technique déjà
établie par le serveur : il a pour fonction de rendre la chaîne
\emph{observable}. Ce qui doit être visible, c'est moins le contenu produit que le
raisonnement qui y conduit — quel état a été perçu par chaque modalité, comment le
désaccord a été tranché, quel patron a été retenu, quelles sources ont ancré la
génération.

\subsection{Interface de visualisation}
\label{subsec:impl-interface}

\paragraph{Une application d'une seule page.} L'interface est une application
\textsc{React} en page unique, sans routeur : la configuration, l'attente et les
résultats occupent des sections successives d'une même vue, ce qui maintient sous
les yeux le lien entre le profil choisi et l'adaptation obtenue. L'adresse du
serveur est lue dans une variable d'environnement, avec repli sur l'hôte local, et
la disponibilité de l'API est vérifiée toutes les trente secondes puis signalée par
un indicateur d'état : un prototype de recherche est fréquemment utilisé pendant que
le serveur redémarre, et distinguer une panne d'API d'une absence de résultat évite
de fausses conclusions lors des démonstrations. Le tableau~\ref{tab:composants-ui}
récapitule les composants et leur rôle.

\begin{table}[H]
\centering
\small
\renewcommand{\arraystretch}{1.15}
\begin{tabular}{|>{\raggedright\arraybackslash}p{4.0cm}|>{\raggedright\arraybackslash}p{8.9cm}|}
\hline
\rowcolor{gray!20}
\textbf{Composant} & \textbf{Rôle dans l'observation de la chaîne} \\
\hline
Sélecteur de scénario & Liste filtrable des cinquante scénarios, avec repère
coloré de l'état dominant et raccourcis vers un scénario type de chaque patron. \\
\hline
Diagramme en barres & Intensité des quatre dimensions du scénario sélectionné,
avec mise en évidence de la dimension dominante. \\
\hline
Diagramme radial & Même information sous forme de profil, plus lisible pour
comparer deux scénarios de mémoire. \\
\hline
Flux de pipeline & Progression en quatre étapes : engagement, décision,
récupération documentaire, contenu. \\
\hline
Panneau d'attente & Messages d'étape pendant l'analyse, dont la durée dépend du
service distant. \\
\hline
Pastille d'état & Codage couleur et symbole constants pour chaque état
d'engagement, réutilisés dans tout l'écran. \\
\hline
Tuteur conversationnel & Dialogue prolongeant l'adaptation, avec activation ou
désactivation de l'ancrage documentaire. \\
\hline
Historique latéral & Dix dernières adaptations, rechargées depuis le serveur au
démarrage. \\
\hline
\end{tabular}
\caption{Composants de l'interface de simulation et fonction de chacun.}
\label{tab:composants-ui}
\end{table}

\paragraph{Des visualisations écrites à la main.} Les deux représentations
graphiques ne reposent sur aucune bibliothèque de tracé : le diagramme en barres est
construit en éléments de mise en page et le diagramme radial est un dessin
vectoriel direct. Pour quatre dimensions bornées à quatre niveaux, une bibliothèque
générique aurait ajouté une dépendance volumineuse sans apporter de fonction utile,
et aurait rendu moins direct le contrôle du codage couleur, lequel doit rester
strictement identique entre les pastilles d'état, les graphiques et les cartes de
fusion.

\subsection{Lecteur de parcours d'apprentissage}
\label{subsec:impl-parcours}

\paragraph{Le parcours comme enchaînement de clips.} Un parcours est, dans le
prototype, une suite ordonnée de clips appartenant à une même séquence
(section~\ref{subsec:impl-segmentation}). Sa lecture s'effectue en enchaînant les
analyses : l'utilisateur sélectionne successivement les clips de la séquence, chaque
analyse produisant une décision et, le cas échéant, un contenu. L'interface
matérialise le parcours par un historique latéral qui conserve les dix dernières
adaptations avec leur état et leur horodatage, et qui est rechargé depuis la base de
sessions au démarrage : une trajectoire simulée survit donc au rafraîchissement de
la page et peut être reprise ou présentée après coup.

\paragraph{Ce que le lecteur ne fait pas.} La lecture est \textbf{pilotée par
l'utilisateur} et non automatisée : l'interface ne comporte ni commande de lecture
continue, ni boutons de clip suivant et précédent, et le référentiel de séquences
n'est pas exposé par le serveur — il reste un artefact hors ligne exploité par les
scripts d'évaluation. Cette limite est cohérente avec le choix, assumé en
section~\ref{subsec:impl-transitions}, d'un point d'entrée au grain du clip ; elle
implique en retour que l'évaluation des trajectoires du chapitre~\ref{chap:validation}
s'appuie sur les séquences rejouées hors ligne, non sur des captures d'écran de
l'interface. L'ajout d'un véritable lecteur de parcours, doté d'une notion de
session identifiée, figure parmi les développements proposés au
chapitre~\ref{chap:conclusion}.

\subsection{Affichage des adaptations}
\label{subsec:impl-affichage}

\paragraph{Rendre la décision lisible avant le contenu.} L'écran de résultats place
la décision avant le texte produit. Un bandeau annonce l'état détecté et son
intensité ; une grille de quatre indicateurs donne le patron retenu, sa famille, le
fait qu'une intervention a bien été déclenchée et le nombre de documents injectés ;
une carte reprend en clair la description de la décision. Ce n'est qu'ensuite
qu'apparaît le contenu, accompagné de son patron, de son objectif et de son nombre
de mots. L'ordre est intentionnel : il empêche de lire le contenu comme un texte
autonome et rappelle qu'il est la conséquence d'une décision documentée.

\paragraph{Montrer la fusion et ses désaccords.} Lorsque l'analyse est
effectivement multimodale, trois cartes juxtaposent l'état perçu par le visage,
l'état perçu par la voix et l'état fusionné. C'est l'élément d'interface le plus
directement lié à la contribution scientifique du travail : il rend visible, cas par
cas, le désaccord entre modalités et la manière dont il a été tranché, là où un
affichage du seul état final laisserait croire à une perception univoque. La liste
des sources documentaires injectées joue le même rôle pour la génération, en
rattachant le texte produit à des références identifiables.

\paragraph{Exploitation du contenu produit.} Deux fonctions closent l'écran : la
copie du contenu dans le presse-papiers et son export en document \textsc{PDF}
titré, portant le concept, l'état, l'intensité et le patron employé. Cette
seconde fonction n'est pas un ornement : elle a servi à constituer les pièces
soumises aux évaluateurs humains du chapitre~\ref{chap:validation}, en garantissant
que le contenu évalué est bien celui produit par le système et qu'il reste
accompagné des métadonnées qui en conditionnent l'appréciation.

%==============================================================================
\section{Génération du corpus de contenu adaptatif}
\label{sec:impl-corpus}

L'artefact final du travail n'est pas seulement un logiciel : c'est un
\emph{corpus} de contenus pédagogiques adaptatifs, produit par le système et
soumis à évaluation. Cette section décrit sa constitution, depuis le choix des
situations couvertes jusqu'à sa documentation.

\subsection{Sélection de 50 séquences représentatives}
\label{subsec:impl-selection-sequences}

\paragraph{Cinquante situations, six patrons.} Les cinquante séquences retenues
sont les clips de dix secondes du référentiel constitué en
section~\ref{subsec:impl-scenarios}. Leur sélection répond à un critère de
couverture et non de représentativité statistique : chacun des six patrons
pédagogiques est présent, les situations problématiques étant sur-représentées pour
que la chaîne de décision soit sollicitée sur l'ensemble de son domaine. Le passage
complet du banc d'essai sur ces cinquante séquences — perception faciale réelle,
perception vocale appariée, fusion, décision — produit une table de décisions
consignant pour chaque scénario l'état retenu, la stratégie et le patron, qui
constitue la base de référence des analyses du chapitre~\ref{chap:validation}.

\paragraph{Des séquences aux trajectoires.} Ces cinquante situations sont en outre
regroupées en dix trajectoires de quatre à sept clips, chacune illustrant une
évolution typique. Une trajectoire n'ajoute aucune donnée : elle ordonne des
scénarios existants pour qu'une dynamique devienne observable. C'est sur ces
trajectoires que les règles temporelles peuvent être exercées, une situation isolée
ne permettant par construction de tester que les règles instantanées.

\subsection{Génération de contenus pour chaque scénario}
\label{subsec:impl-generation-corpus}

\paragraph{Volume et couverture.} Le corpus a été produit en appliquant la chaîne
complète à chaque scénario, pour plusieurs patrons compatibles avec son état, de
manière à couvrir la bibliothèque au-delà des seuls patrons que la cascade
sélectionne spontanément. Il compte environ \textbf{220 contenus}, chacun
enregistré avec l'identifiant du clip, son patron pédagogique, le patron de prompt
employé, l'émotion et son intensité, le texte produit et son nombre de mots.
Dix-neuf des vingt patrons de prompts sont représentés, contre onze lors du seul
passage automatique sur les cinquante scénarios : le corpus documente donc
davantage de situations pédagogiques que le référentiel n'en produit
spontanément. Le patron manquant, dédié à un exercice comparatif, n'a été atteint
par aucune combinaison de scores présente dans l'échantillon — observation
cohérente avec la concentration de la cascade relevée en
section~\ref{subsec:impl-selection}.

\paragraph{Caractéristiques mesurées.} Les indicateurs de forme calculés sur
l'ensemble du corpus situent les contenus autour de \textbf{225 mots} en moyenne,
dans une plage allant de 115 à 434 mots, avec un taux de répétition lexicale moyen
proche de 0{,}49. La dispersion des longueurs suit les longueurs cibles des patrons :
les patrons d'encouragement produisent les textes les plus courts, ceux d'exercice
et d'application les plus longs. Cette conformité relative n'est pas garantie par
construction — aucun mécanisme ne tronque les sorties
(section~\ref{subsec:impl-postproc}) — et constitue donc un résultat en soi,
discuté au chapitre~\ref{chap:validation}.

\paragraph{Régénération ciblée plutôt que globale.} Le corpus n'a pas été produit
en une passe. Un premier audit a signalé un nombre important de contenus
incohérents, principalement des associations douteuses entre émotion enregistrée et
famille de patron. Une passe de correction des seules métadonnées, recalculant la
décision sans rappeler le modèle, a ramené la proportion de contenus conformes à
plus de neuf sur dix ; une passe de régénération, limitée aux lignes encore
signalées et sauvegardant le corpus avant écriture, a porté cette proportion à
\textbf{97,7 \%}, les anomalies résiduelles ne concernant plus que cinq contenus.
Deux propriétés de l'outillage ont rendu cette progression possible : la
régénération est \emph{sélective}, elle ne retouche que les entrées signalées par
l'audit et laisse les autres intactes ; et elle est \emph{temporisée}, un délai
séparant deux appels au service distant, sans quoi la limitation de débit
interromprait le traitement. Cette traçabilité du corpus — état initial, correction
appliquée, état final — fait partie du livrable au même titre que les contenus.

\subsection{Documentation du corpus}
\label{subsec:impl-doc-corpus}

\paragraph{Artefacts livrés.} Le corpus n'est pas un fichier isolé mais un ensemble
d'artefacts articulés (tableau~\ref{tab:artefacts-corpus}) : les contenus
eux-mêmes, le rapport d'audit ligne par ligne, les indicateurs de forme,
l'échantillon stratifié soumis aux évaluateurs et les fiches d'évaluation vierges.
Chaque fichier est produit par un script identifié, à un emplacement fixe, ce qui
permet de reconstituer la chaîne complète depuis les scénarios jusqu'aux fiches.

\begin{table}[H]
\centering
\small
\renewcommand{\arraystretch}{1.2}
\begin{tabular}{|>{\raggedright\arraybackslash}p{4.4cm}|>{\raggedright\arraybackslash}p{6.5cm}|>{\raggedright\arraybackslash}p{2.0cm}|}
\hline
\rowcolor{gray!20}
\textbf{Artefact} & \textbf{Contenu} & \textbf{Volume} \\
\hline
Corpus de contenus & Texte généré, patron, émotion, intensité, nombre de mots &
$\approx$ 220 \\
\hline
Rapport d'audit & Anomalies détectées et score de qualité par contenu &
$\approx$ 220 \\
\hline
Indicateurs de forme & Longueur, phrases, lisibilité approchée, répétition &
$\approx$ 220 \\
\hline
Échantillon d'évaluation & Tirage stratifié par famille de patron & 58 \\
\hline
Fiches d'évaluation & Extrait de contenu et grille de notation à remplir & 58 \\
\hline
Table de décisions & État, stratégie et patron retenus par scénario & 50 \\
\hline
\end{tabular}
\caption{Artefacts constituant le corpus de contenu adaptatif et son appareil de
validation.}
\label{tab:artefacts-corpus}
\end{table}

\paragraph{Documentation méthodologique.} Deux documents accompagnent ces fichiers.
Le premier décrit la procédure de régénération : commandes exactes de l'audit, de
la simulation à blanc, de la correction des métadonnées puis de la régénération
complète, prérequis d'accès au service de génération, et proportions de conformité
obtenues à chaque étape. Le second fixe le protocole d'évaluation : constitution de
l'échantillon à quinze contenus par famille de patron, maintien d'un concept
identique pour rendre les contenus comparables, critères de notation soumis aux
enseignants et articulation avec les indicateurs automatiques. Ce protocole est
détaillé et mis en œuvre au chapitre~\ref{chap:validation}.

\paragraph{Reproductibilité de la chaîne.} Le tirage de l'échantillon d'évaluation
est lui-même initialisé par une graine fixe, et les fiches soumises aux évaluateurs
sont générées par script à partir du corpus, non recopiées manuellement. L'écart
entre l'effectif visé — soixante contenus, quinze par famille — et l'effectif
obtenu de cinquante-huit provient de la couverture inégale des familles de patrons
dans le corpus ; il est mentionné ici pour que les effectifs rapportés au
chapitre~\ref{chap:validation} soient interprétés à leur juste valeur.

\paragraph{Bilan de l'implémentation.} Au terme de ce chapitre, l'artefact est
complet : deux corpus préparés et documentés, une chaîne de perception dont les
variantes sont commutables par configuration, un moteur de décision dont chaque
règle est nommée et testée, un module de génération ancré documentairement, une
interface qui expose le raisonnement autant que son résultat, et un corpus de
contenus accompagné de son appareil de validation. Trois limites ont été énoncées
en chemin plutôt que dissimulées : l'absence de mémoire entre décisions
successives, la lecture de parcours pilotée manuellement et l'absence de
post-traitement des contenus générés. Le chapitre~\ref{chap:validation} soumet cet
artefact à l'épreuve des faits, en mesurant successivement la perception, la
décision et la qualité pédagogique du corpus produit.
